\begin{table}[htbp]
\centering\footnotesize
\caption{Element loadings on the laboratory-XRF elemental axis (PC1 after within-campaign standardization).}
\label{tab:xrf_loadings}
\setlength{\tabcolsep}{4pt}
\begin{tabular}{lrrrrrrrrrrr}
\toprule
Element & Ca & Mg & Na & S & Cl & Fe & Ti & K & Al & P & Si \\
\midrule
Loading & $0.41$ & $0.41$ & $0.37$ & $0.35$ & $0.34$ & $0.28$ & $0.28$
        & $0.17$ & $0.16$ & $-0.01$ & $-0.27$ \\
\bottomrule
\end{tabular}
\end{table}

\begin{table}[htbp]
\centering\footnotesize
\caption{Genera whose site-level CLR values correlate with the site-mean laboratory XRF axis (Spearman, $60$ sites). $10$ with the largest $|\rho|$ in each direction are shown. Positive $\rho$ indicates higher abundance towards the Ca, Mg, Na, S and Cl end of the axis, negative $\rho$ towards the Si end.}
\label{tab:xrf-genera}
\begin{tabular}{llrl}
\toprule
Genus & Phylum & $\rho$ & $q$ \\
\midrule
\multicolumn{4}{l}{\textit{Increasing along the axis}} \\
\textit{Polygonibacillus}     & Bacillota      & $0.81$ & $1.8\times10^{-12}$ \\
\textit{Escherichia-Shigella} & Pseudomonadota & $0.74$ & $6.2\times10^{-10}$ \\
\textit{Sediminibacillus}     & Bacillota      & $0.74$ & $6.2\times10^{-10}$ \\
\textit{Aquibacillus}         & Bacillota      & $0.74$ & $6.2\times10^{-10}$ \\
\textit{Halalkalibacter}      & Bacillota      & $0.73$ & $1.2\times10^{-9}$ \\
\textit{Gracilibacillus}      & Bacillota      & $0.72$ & $3.9\times10^{-9}$ \\
\textit{Enterococcus}         & Bacillota      & $0.71$ & $5.4\times10^{-9}$ \\
\textit{Aeromonas}            & Pseudomonadota & $0.69$ & $1.7\times10^{-8}$ \\
\textit{Pseudalkalibacillus}  & Bacillota      & $0.69$ & $2.5\times10^{-8}$ \\
\textit{Crossiella}           & Actinomycetota & $0.65$ & $2.3\times10^{-7}$ \\
\addlinespace
\multicolumn{4}{l}{\textit{Decreasing along the axis}} \\
\textit{Steroidobacter}       & Pseudomonadota  & $-0.71$ & $5.4\times10^{-9}$ \\
\textit{Caldilinea}           & Chloroflexota   & $-0.67$ & $9.3\times10^{-8}$ \\
\textit{Phenylobacterium}     & Pseudomonadota  & $-0.65$ & $2.7\times10^{-7}$ \\
\textit{Paraconexibacter}     & Actinomycetota  & $-0.65$ & $3.0\times10^{-7}$ \\
\textit{Streptomyces}         & Actinomycetota  & $-0.65$ & $3.1\times10^{-7}$ \\
\textit{Roseisolibacter}      & Gemmatimonadota & $-0.63$ & $7.1\times10^{-7}$ \\
\textit{Gemmatimonas}         & Gemmatimonadota & $-0.63$ & $7.2\times10^{-7}$ \\
\textit{Pirellula}            & Planctomycetota & $-0.58$ & $8.9\times10^{-6}$ \\
\textit{Blastocatella}        & Acidobacteriota & $-0.58$ & $8.9\times10^{-6}$ \\
\textit{Dongia}               & Pseudomonadota  & $-0.58$ & $8.9\times10^{-6}$ \\
\bottomrule
\end{tabular}
\end{table}